% !TeX root = ../../Book4.tex
% 纲要标题：### 20.3 倾斜对象
% 纲要位置：计划纲要.md，第 3039 行。

\section{倾斜对象}
\label{b4:sec:2003}


% 纲要内容（仅作写作计划，尚非教材正文）：
%
% * 消失条件与生成条件
% * 经典有限维代数例子
% * 导出 Morita 理论的基本版本
%

\subsection{消失条件与生成条件}
\label{b4:subsec:2003:01}

第二卷定理16.3.3已经证明，普通 Morita 理论用有限生成投射生成元代替正则模。导出范畴允许把这个替代对象选成复形，不过此时必须同时控制它与自身各次移位之间的态射。如果这些高次态射不消失，端同态信息便应由一个 DG 代数承担，不能只留下普通端同态环。

\begin{definition}\label{b4:def:tilting-complex}
设 \(k\) 为域，\(R\) 为含幺结合 \(k\)-代数。若 \(T\in\operatorname{Perf}(R)\) 满足
\[
 \operatorname{Hom}_{D(R\text{-}\mathrm{Mod})}(T,T[n])=0
 \quad(n\ne0),\qquad
 \operatorname{thick}(T)=\operatorname{Perf}(R),
\]
称 \(T\) 为 \(R\) 上的\term[qingxiefuxing]{倾斜复形}[tilting complex]，也称所给三角范畴中的倾斜对象。这里使用的是包括正负所有非零次数的消失条件。
\end{definition}

\(R\) 本身是倾斜复形：其自身高次态射为零，厚生成条件由\cref{b4:thm:perfect-compact-thick}给出。两个条件不能相互代替。若 \(R=k\times k\)，直和项 \(k\times0\) 没有非零高次自身态射，却生成不了 \(0\times k\)。若 \(R=k\)，对象 \(k\oplus k[1]\) 生成全部完美复形，但它到一次移位的态射含有从原来 \(k[1]\) 到移位后 \(k[1]\) 的恒等映射，消失条件失败。

\ProseParagraphBreak
这里的倾斜对象与\chapref{b4:ch:19}的扭对倾斜使用不同数据。后者从一个心及其中的扭对构造另一个心；本节则要求一个完美对象具有自身态射消失和厚生成性质。构造出新心并不自动产生这样的对象，更不自动给出两个模导出范畴之间的等价。

\begin{proposition}\label{b4:prop:tilting-module-criterion}
设 \(k\) 为域，\(R\) 为 \(k\)-代数，\(T\) 为左 \(R\)-模。假设 \(T\) 有有限长有限生成投射消解，\(\operatorname{Ext}^n_R(T,T)=0\) 对所有 \(n>0\) 成立，并且存在正合列
\[
 0\longrightarrow R\longrightarrow T_0\longrightarrow T_1
 \longrightarrow\cdots\longrightarrow T_m\longrightarrow0,
\]
其中每个 \(T_i\) 是某个有限直和 \(T^{r_i}\) 的直和项。则 \(T\) 集中在零次是倾斜复形。
\end{proposition}
\begin{proof}
有限投射消解给出完美性。对 \(n>0\)，自身导出态射由\cref{b4:thm:derived-ext-identification}识别为所给 Ext 群；对 \(n<0\)，取位于非正次数的投射消解 \(P\to T\)，则 \(\underline{\operatorname{Hom}}_R(P,T)\) 没有负次项，故负次态射为零。

把所给长正合列按像分成短正合列。从末项起逆向使用相应好三角，可知 \(R\in\operatorname{thick}(T)\)。由\cref{b4:thm:perfect-compact-thick}，\(\operatorname{Perf}(R)=\operatorname{thick}(R)\subseteq\operatorname{thick}(T)\)。反向包含由完美子范畴的厚性和 \(T\) 完美得到。
\end{proof}

\subsection{经典有限维代数例子}
\label{b4:subsec:2003:02}

下面用三个向量空间及两个线性映射计算一个倾斜对象，暂不需要第五卷的一般箭图表示论。设 \(k\) 为域，\(R\) 是箭图 \(1\to2\to3\) 的路代数：以三条长度零的路、两条箭头及它们的复合为 \(k\)-基，乘法按可接续路径复合，不可接续时取零。具体记箭头为 \(\alpha:1\to2\)、\(\beta:2\to3\)，约定 \(\alpha=e_2\alpha e_1\)、\(\beta=e_3\beta e_2\)，所以 \(\beta\alpha\) 是从一到三的路。有限维左 \(R\)-模等价于图式 \(V_1\to V_2\to V_3\)；三个顶点幂等元给出三个分量，箭头作用给出两个线性映射。模同态是使这两个方块交换的三元线性映射。

\ProseParagraphBreak
写出三个投射模以及中间顶点的简单模：
\[
\begin{aligned}
 P_1&=(k\xrightarrow{1}k\xrightarrow{1}k),&
 P_2&=(0\longrightarrow k\xrightarrow{1}k),\\
 P_3&=(0\longrightarrow0\longrightarrow k),&
 S_2&=(0\longrightarrow k\longrightarrow0).
\end{aligned}
\]
这里 \(P_i=Re_i\)，并且 \(\operatorname{Hom}_R(P_i,V)\cong V_i\)：态射由顶点 \(i\) 处的向量决定，其余分量由箭头推出。因此它们确为投射模，且 \(R=P_1\oplus P_2\oplus P_3\)。

\begin{example}\label{b4:exm:a3-tilting-object}
在上述路代数 \(R\) 上，令 \(T=P_1\oplus P_2\oplus S_2\)。则 \(T\) 是倾斜复形，而
\[
 B=\operatorname{End}_R(T)^{\mathrm{op}}
\]
同构于路代数 \(k(1\to2\leftarrow3)\)。
\end{example}
\begin{proof}
有短正合列
\begin{equation}\label{b4:eq:a3-simple-resolution}
 0\longrightarrow P_3\xrightarrow{i}P_2
 \xrightarrow{q}S_2\longrightarrow0.
\end{equation}
所以 \(T\) 有长度至多一的有限投射消解。对任意表示 \(V\)，这份消解给出
\[
 \operatorname{Ext}^1_R(S_2,V)
 =\operatorname{coker}(V_2\longrightarrow V_3),
 \qquad \operatorname{Ext}^{n}_R(S_2,V)=0\quad(n\ge2).
\]
对 \(V=P_1,P_2,S_2\)，\(V_2\to V_3\) 都满；因此 \(\operatorname{Ext}^n_R(T,T)=0\) 对所有 \(n>0\) 成立。负次态射为零同样由非正次投射消解计算。\eqref{b4:eq:a3-simple-resolution}的三角把 \(P_3\) 从 \(P_2,S_2\) 构造出来，故 \(R\in\operatorname{thick}(T)\)，由完美性即得厚生成。

端同态的计算也可以逐项完成。三个直和项的端同态环均为 \(k\)。不同项之间仅有
\[
 \operatorname{Hom}_R(P_2,P_1)=k j,
 \qquad \operatorname{Hom}_R(P_2,S_2)=k q
\]
非零：从 \(P_1\) 出发的态射由目标顶点一的分量决定，故另外两项为零；从 \(S_2\) 到 \(P_1\) 或 \(P_2\) 的态射则由第二条箭头的交换条件迫为零。相对于 \(T\) 的三个直和项，端同态代数因此为
\[
 E=\left\{\left.
 \begin{pmatrix}a&b&0\\0&c&0\\0&d&e\end{pmatrix}
 \ \right|\ a,b,c,d,e\in k\right\}.
\]
其中两条非恒等箭头从第二个直和项发出，彼此不能复合。取反代数后，两条箭头都指向第二个顶点，正好得到 \(B=k(1\to2\leftarrow3)\)。
\end{proof}

这个例子不是普通 Morita 等价换了一种写法。\(R\) 的三个不可分解投射模的合成长度为 \(3,2,1\)；\(B\) 的三个不可分解投射模分别支持在 \(1\to2\)、顶点二、\(3\to2\)，合成长度为 \(2,1,2\)。这些类型已经穷尽：两代数模去箭头生成的幂零理想均为 \(k^3\)，三个顶点给出全部简单模，第二卷定理15.3.10的简单模与不可分解投射模对应给出所述清单。模范畴等价保持简单对象、投射性、不可分解性和合成列，因而保持这组长度；两组数不同，所以两者的模范畴不等价。下一小节的定理却会给出其导出范畴的等价。

\subsection{导出 Morita 理论的基本版本}
\label{b4:subsec:2003:03}

令 \(T\) 为倾斜复形，选择有界有限生成投射模型 \(P\)，并令
\[
 E=\underline{\operatorname{End}}_R(P),\qquad
 B=\operatorname{End}_{D(R\text{-}\mathrm{Mod})}(T)^{\mathrm{op}}.
\]
倾斜条件恰好说 \(H^n(E)=0\) 对 \(n\ne0\) 成立，而 \(H^0(E)=B^{\mathrm{op}}\)。这一步把一个涉及所有移位的条件变为端同态 DG 代数的上同调消失。

\begin{proposition}\label{b4:prop:dg-degree-zero-formality}
设 \(k\) 为交换环，\(E\) 为 \(k\) 上的 DG 代数，且 \(H^n(E)=0\) 对所有 \(n\ne0\) 成立。令
\[
 (\tau_{\le0}E)^n=
 \begin{cases}
 E^n,&n<0,\\
 \ker(d:E^0\to E^1),&n=0,\\
 0,&n>0.
 \end{cases}
\]
原微分和乘法使 \(\tau_{\le0}E\) 成为 DG 子代数，而且包含映射与零次取商给出 DG 代数拟同构
\[
 E\ \longleftarrow\ \tau_{\le0}E\ \longrightarrow\ H^0(E).
\]
\end{proposition}
\begin{proof}
非正次数在乘法下封闭；两零次余圈之积仍为余圈。来自 \(E^{-1}\) 的微分落在零次余圈中，故这是子复形。到 \(H^0(E)\) 的映射在负次取零，在零次把余圈取模余边。负次元素的乘积不会产生正次或零次项；两个零次余圈的乘法与取商相容，所以它是代数映射。包含映射在非正次上同调上为同构，而 \(E\) 正次上同调为零；取商映射在零次为同构，两边其他次数上同调均为零。因此两映射都是拟同构。
\end{proof}

注意这里得到的是通过第三个 DG 代数连接的两个拟同构，未必有直接代数映射 \(H^0(E)\to E\)。任选每个同伦类的复形映射代表，通常不能同时保持乘法。正是这一点阻止了把普通 Morita 理论中的论证逐字照搬。

\begin{theorem}[导出 Morita 定理的倾斜版本]\label{b4:thm:derived-morita-tilting}
设 \(k\) 为域，\(R\) 为含幺结合 \(k\)-代数，\(T\in\operatorname{Perf}(R)\) 为倾斜复形，并令
\[
 B=\operatorname{End}_{D(R\text{-}\mathrm{Mod})}(T)^{\mathrm{op}}.
\]
则存在 \(R\)-\(B\) 双模复形 \(X\)，作为左 \(R\)-模复形时在导出范畴中同构于 \(T\)，使
\[
 X\otimes_B^{\mathbf L}-:
 D(B\text{-}\mathrm{Mod})\ \longrightarrow\ D(R\text{-}\mathrm{Mod})
\]
成为 \(k\)-线性三角等价，且把 \(B\) 送到 \(T\)。此等价限制为 \(\operatorname{Perf}(B)\simeq\operatorname{Perf}(R)\)。
\end{theorem}

这是本章外引的\term[daochuMoritalilun]{导出 Morita 定理}[derived Morita theorem]基本版本。Keller 的\cite[Section 4, Theorem (Rickard)]{Keller1998AbelianDerived}给出普通代数版本；经 DG 端同态代数及截断实现等价的方法见\cite[Sections 8.2 and 9.1--9.2]{Keller1994DerivingDG}。这些文献主要用右模；本定理改用左模，故端同态环必须取反。这里选择域为底环，使普通代数的底模平坦性自动成立。

\ProseParagraphBreak
本书已经给出了定理中的三项可检查数据：有界有限生成投射模型、自身移位态射消失、有限的厚生成过程。真正没有在本章重证的是从同伦意义的 \(B\) 作用构造严格双模复形 \(X\)，以及在无界 DG 模导出范畴中证明张量等价。其构造先经 \(E^{\mathrm{op}}\)，再经\cref{b4:prop:dg-degree-zero-formality}的两个拟同构。因此不能把“\(B\) 在导出范畴中作用于 \(T\)”直接当作“所选 \(P\) 已逐项成为右 \(B\)-模”。完美子范畴上的限制则无需另添假设：三角等价保持直和项，而 \(B\) 的像为 \(T\)，所以把 \(\operatorname{thick}(B)\) 送到 \(\operatorname{thick}(T)\)。

\ProseParagraphBreak
把定理用于\cref{b4:exm:a3-tilting-object}，便得到 \(k(1\to2\to3)\) 与 \(k(1\to2\leftarrow3)\) 的导出等价。虽然它们的模范畴不同，允许使用复形后，原来的投射生成元可以由 \(P_1\oplus P_2\oplus S_2\) 代替；\eqref{b4:eq:a3-simple-resolution}正是恢复缺少的投射模 \(P_3\) 的计算。
